% !TeX program = lualatex
%======================================================================
% DOCUMENT SETUP -- compile with recent LuaLaTeX
%======================================================================
% Metadata must appear before \documentclass.
\DocumentMetadata{
  tagging=on,
  tagging-setup={math/setup=mathml-SE},
  pdfstandard=ua-2,
  lang=en-US
}
\documentclass[10pt,letterpaper]{article}
\usepackage{notes}
% \usepackage[widemargin]{notes}

%----------------------------------------------------------------------
% PDF PROPERTIES -- title for this lecture; author/course from notes.sty
%----------------------------------------------------------------------
\hypersetup{
  pdflang={en-US},
  pdftitle={\class{} - 2.6 Method of Undetermined Coefficients - Part 1},
  pdfauthor={\instructor},
  pdfsubject={\coursename}
}

% Standalone lecture: compile this file from the repository root.
% Keep notes.sty on the TeX search path. The original 2.6-mouc.tex is unchanged.
% Timing comments describe a suggested 9:00--9:50 class; optional extensions
% remain in the notes for independent reading or a longer class.
\begin{document}
\lesson[Logan, Differential Equations, 3rd edition\enspace\textbullet\enspace Section 2.3.1]
  {2.6, Part 1}{Building the Method}
\renewcommand{\lessontitle}{Undetermined Coefficients, Part 1}
\section{Method of Undetermined Coefficients}

\subsection{Homogeneous vs. Nonhomogeneous}

% SIDE NOTE: 9:00 | Two Pieces
% a) A nonzero forcing term changes the problem.
% b) Separate the homogeneous family from one particular solution.
% c) Keep the integrating-factor review brief.
Throughout these two lectures, $a,b,c$ are real constants and $a\ne0$.
We are studying linear equations with constant coefficients.

So far, we have been solving \textbf{homogeneous problems} (HP):
\begin{equation}
ax''+bx'+cx=0. \tag{HP}
\end{equation}
Homogeneous means that the right-hand side is zero. A
\textbf{nonhomogeneous problem} (NHP) has the form
\begin{equation}
ax''+bx'+cx=f(t), \tag{NHP}
\end{equation}
where $f(t)$ is not identically zero. The function $f(t)$ is
called the \textbf{forcing term} or \textbf{source term}.
For a spring-mass system, it represents an external force
acting on the mass. Our goal is to solve nonhomogeneous equations
for certain types of forcing terms. We begin with a
first-order example.

%-------------------------------------------------------------
\subsection{First-Order NHP}

Consider the equation
\[
x'+2x=9e^t.
\]
We already know how to solve this using an integrating factor.
With $\mu(t)=e^{2t}$, we obtain
\begin{align*}
x(t)
&=e^{-2t}\left(\int 9e^{3t}\,dt+C\right)\\
&=e^{-2t}\left(3e^{3t}+C\right)\\
&=Ce^{-2t}+3e^t.
\end{align*}
Notice that the solution has two pieces:
\[
x(t)=
\underbrace{Ce^{-2t}}_{\text{solves }x'+2x=0}
+
\underbrace{3e^t}_{\text{solves }x'+2x=9e^t}.
\]
We will use this structure to develop a new method.

%-------------------------------------------------------------
\subsection{New Method}

We solve the same equation by finding these two pieces
separately:
\[
x(t)=x_h(t)+x_p(t).
\]
Here, $x_h$ is the \textbf{general homogeneous solution},
and $x_p$ is a \textbf{particular solution}: any one
solution of the nonhomogeneous equation.

\par\medskip
\textbf{Step 1. Find the homogeneous solution.}
Temporarily set the forcing term to zero:
\[
x'+2x=0.
\]
The general solution is
\[
x_h(t)=Ce^{-2t}.
\]

\par\medskip
\textbf{Step 2. Choose a trial function.}
Since the forcing term is $9e^t$, we guess a form for the particular solution
\[
x_p(t)=Ae^t,
\]
where $A$ is an unknown coefficient. Differentiating $e^t$
preserves its form, so this is a reasonable guess.

\par\medskip
\textbf{Step 3. Determine the coefficient.}
Substitute the trial function into the nonhomogeneous equation:
\begin{align*}
x_p'+2x_p&=9e^t\\
Ae^t+2Ae^t&=9e^t\\
3Ae^t&=9e^t.
\end{align*}
Thus, $A=3$ and $x_p(t)=3e^t$.

\par\medskip
\textbf{Step 4. Combine the two pieces.}
The general solution is
\[
\boxed{x(t)=x_h(t)+x_p(t)=Ce^{-2t}+3e^t.}
\]
This agrees with the integrating-factor solution.
The same approach also works for second-order equations.

%-------------------------------------------------------------
\subsection{Why the Sum Works}
Why does the summation 
$$
    x=x_h+x_p
$$
work?
Suppose $x_h$ solves the homogeneous equation and $x_p$
solves the nonhomogeneous equation. Substituting their sum gives
\begin{align*}
ax''+bx'+cx&=f(t)\\
a(x_h+x_p)''+b(x_h+x_p)'+c(x_h+x_p)&=f(t)\\
(ax_h''+bx_h'+cx_h)+(ax_p''+bx_p'+cx_p)&=f(t)\\
(0)+(f(t))&=f(t)\\
f&=f \; \checkmark
\end{align*}
Therefore, $x_h+x_p$ solves the nonhomogeneous equation. There is also a reason this gives \emph{every} solution. If $x$ is any
solution of the nonhomogeneous equation, subtract the particular solution:
\[
a(x-x_p)''+b(x-x_p)'+c(x-x_p)=f(t)-f(t)=0.
\]
Thus $x-x_p$ belongs to the homogeneous family, and $x=x_h+x_p$.

A particular solution is not unique. For example, both $3e^t$ and
$3e^t+5e^{-2t}$ solve $x'+2x=9e^t$. We choose the simplest convenient
one; the homogeneous term already accounts for all the remaining freedom.

%-------------------------------------------------------------
\subsection{Method of Undetermined Coefficients}

The \textbf{method of undetermined coefficients} finds a
particular solution by choosing a trial function with unknown
coefficients and solving for them. To solve
\[
ax''+bx'+cx=f(t),
\]
use the following steps:
\begin{enumerate}
    \item Find the general homogeneous solution $x_h$ by solving
    \[
    ax''+bx'+cx=0.
    \]

    \item Choose a trial function for $x_p$ based on the form
    of $f(t)$. 

    \item Substitute $x_p$ into the nonhomogeneous equation
    and solve for its unknown coefficients.

    \item Write the general solution
    \[
    x(t)=x_h(t)+x_p(t).
    \]

    \item If initial conditions are given, apply them to the
    complete solution to determine the arbitrary constants.
\end{enumerate}

We already know how to find $x_h$ using the characteristic
equation. The new task is choosing a suitable trial function
for $x_p$, which we discuss next. 




The coefficients $A,B,\ldots$ in a trial are \textbf{undetermined coefficients}:
substitution into the DE fixes their values. The constants $c_1,c_2$ in $x_h$
are \textbf{arbitrary constants}: initial conditions determine their values.
Do not use the initial conditions to choose $A,B,\ldots$.

% TRANSITION (blue box): We know how to find x_h. How do we choose x_p?
\section{The Trial Function}

\subsection{When the Method Applies}

The method of undetermined coefficients applies when the forcing
term $f(t)$ is built from:
\begin{itemize}
    \item Polynomials,
    \item Exponentials $e^{\beta t}$,
    \item Waves $\sin(\omega t)$ and $\cos(\omega t)$,
    \item Sums and products of these functions.
\end{itemize}
For example, we can use the method to solve
\[
x''+x'+x=3e^{2t}\sin t+t^2,
\]
but it does not apply to
\[
x''+x'+x=\ln t.
\]
The key feature of the allowed functions is that repeated
differentiation produces the same types of terms. This lets
us choose a trial function with finitely many unknown coefficients.

%-------------------------------------------------------------

% SIDE NOTE: 9:10 | A Complete Second-Order Example
% a) Find x_h using the characteristic equation.
% b) Choose a full quadratic and match all three powers.
% c) Explain why the missing linear term still needs a coefficient.
\subsection{A Polynomial Forcing Term}
\begin{example}{
Find the general solution to
\[
x''-3x'+2x=2t^2-3.
\]
}
\textbf{Step 1. Find the homogeneous solution.}
First, solve the homogeneous equation. The characteristic equation is
\[
\lambda^2-3\lambda+2=(\lambda-1)(\lambda-2)=0,
\]
so
\[
x_h(t)=c_1e^t+c_2e^{2t}.
\]

\textbf{Step 2. Choose and differentiate the trial.}
Since the forcing term is a quadratic polynomial, use
\[
x_p(t)=At^2+Bt+C.
\]
We include the linear term even though it does not appear
in the forcing term. Differentiating gives
\[
x_p'=2At+B,
\qquad
x_p''=2A.
\]
\textbf{Step 3. Match coefficients.}
Substitute into the nonhomogeneous equation and collect
like powers of $t$:
\begin{align*}
2A-3(2At+B)+2(At^2+Bt+C)&=2t^2-3,\\
(2A)t^2+(2B-6A)t+(2A-3B+2C)&=2t^2-3.
\end{align*}
For this equality to hold for every $t$, the coefficients
of corresponding powers must match:
\[
\begin{dcases}
2A=2,\\
2B-6A=0,\\
2A-3B+2C=-3.
\end{dcases}
\]
Solving in order gives
\[
A=1,\qquad B=3,\qquad C=2.
\]
\textbf{Step 4. Assemble the general solution.}
Therefore, $x_p(t)=t^2+3t+2$, and the general solution is
\[
\boxed{x(t)=c_1e^t+c_2e^{2t}+t^2+3t+2.}
\]
\end{example}


\subsection{Why Include the Missing Powers?}
In the previous example, the forcing has no linear term. Why include $Bt$?
If we had tried only $At^2+C$, substitution would give
\[
x_p''-3x_p'+2x_p=2At^2-6At+(2A+2C).
\]
Matching $2t^2$ requires $A=1$, but then the unwanted term $-6t$ remains.
The extra term $Bt$ gives us a way to cancel it. A coefficient may turn
out to be zero, but we should not set it to zero before substitution.

Matching coefficients means matching an identity for \emph{every} $t$,
not just making the equation work at one chosen time.
\subsection{Table of Trial Functions}

The table below gives the basic trial forms. Greek letters
represent known constants from the forcing term; capital
Latin letters represent unknown coefficients.

\begin{table}[H]
\centering
\renewcommand{\arraystretch}{1.6}
\begin{tabular}{|c|c|c|}
\hline
\textbf{Type}
& \textbf{Forcing term $f(t)$}
& \textbf{Trial function $x_p(t)$}\\
\hline
Constant
& $\alpha$
& $A$\\
\hline
Exponential
& $\alpha e^{\beta t}$
& $Ae^{\beta t}$\\
\hline
Polynomial
& $\alpha_2t^2+\alpha_1t+\alpha_0$
& $At^2+Bt+C$\\
\hline
Wave
& \shortstack{$\alpha\sin(\omega t)$ or $\alpha\cos(\omega t)$}
& $A\cos(\omega t)+B\sin(\omega t)$\\
\hline
\end{tabular}
\caption{Basic trial forms suggested by the forcing term.}
\end{table}

The polynomial row extends to any degree: for a polynomial
of degree $n$, use a trial polynomial of degree $n$ with an
unknown coefficient for every power of $t$, including the constant. 
These are \textbf{initial trial forms}. Substitution will tell us whether
the proposed form can produce the required forcing term.

%-------------------------------------------------------------
\begin{example}{
Choose an initial trial function from the table for each forcing term.
Do not solve for its coefficients.
\begin{enumerate}
    \item $f(t)=6$
    \item $f(t)=-3e^{4t}$
    \item $f(t)=3t^2$
    \item $f(t)=12\sin 3t$
\end{enumerate}
}
\begin{enumerate}
    \item For a constant forcing term, try
    \[
    x_p(t)=A.
    \]

    \item For an exponential forcing term, keep the
    same exponent:
    \[
    x_p(t)=Ae^{4t}.
    \]

    \item For a quadratic forcing term, include every
    power through degree two:
    \[
    x_p(t)=At^2+Bt+C.
    \]

    \item For a sinusoidal forcing term, include both
    sine and cosine with the same frequency:
    \[
    x_p(t)=A\cos 3t+B\sin 3t.
    \]
\end{enumerate}
\end{example}

The trial function may need terms that are absent from $f(t)$.
Our goal is to include the types of terms produced by
differentiation. For example, if $f(t)=3t^2$, we use
\[
x_p(t)=At^2+Bt+C,
\]
even though the forcing term has no linear or constant term.
Including all three coefficients gives us the freedom to
match the quadratic term and cancel any unwanted lower-degree
terms after substitution. The same goes for waves. We need both 
the sine and cosine because their derivatives switch between the two. 


% SIDE NOTE: 9:23 | Why Both Sine and Cosine?
% a) Differentiation switches sine and cosine.
% b) Include both, then let substitution determine the coefficients.
% c) The trial keeps the forcing frequency.
\subsection{A Sinusoidal Forcing Term}
\begin{example}{Find the general solution to
\[
x''+2x'+2x=5\cos t.
\]}
\textbf{Step 1. Find the homogeneous solution.}
The characteristic equation is $\lambda^2+2\lambda+2=0$, so
\[
x_h=e^{-t}(c_1\cos t+c_2\sin t).
\]
\textbf{Step 2. Choose the trial.}
Try $x_p=A\cos t+B\sin t$. Then
\[
x_p'=-A\sin t+B\cos t,\qquad x_p''=-A\cos t-B\sin t.
\]
\textbf{Step 3. Match the sine and cosine coefficients.}
Substitution gives
\[
(A+2B)\cos t+(B-2A)\sin t=5\cos t+0\sin t.
\]
Consequently,
\[
A+2B=5,\qquad B-2A=0,
\]
so $A=1$ and $B=2$. Therefore,
\[
\boxed{x=e^{-t}(c_1\cos t+c_2\sin t)+\cos t+2\sin t.}
\]
\end{example}
Trying only $A\cos t$ would produce an unwanted sine term from $2x_p'$.
Including $B\sin t$ lets us cancel that term. Notice also that the
frequency is fixed by the forcing: we determine the amplitudes, not a new frequency.

% SIDE NOTE: 9:32 | Sums and Products
% a) Add trial groups for sums.
% b) Combine factors for products, keeping coefficients independent.
% c) Do not solve a large coefficient system here.
\section{Rules for Trial Functions}

The basic trial forms extend to more complicated forcing terms.
We begin with two rules: the sum rule and the product rule.

%-------------------------------------------------------------
\subsection{Sum Rule}

If the forcing term is a sum, add the trial forms for its
individual terms. Use separate unknown coefficients for
each part.

For example, if
\[
x''+2x=t+\sin 3t,
\]
then the initial trial function is
\[
x_p(t)=
\underbrace{At+B}_{\text{for }t}
+
\underbrace{C\cos 3t+D\sin 3t}_{\text{for }\sin 3t}.
\]
We can substitute the entire trial function at once or
find a particular solution for each forcing term separately
and add the results. For instance, you may substitute 
$$
At+B
$$
into $x''+2x=t$ and solve for $A$ and $B$. Then you can 
substitute
$$
C\cos3t+D\sin3t
$$
into $x''+2x=\sin3t$ to find $C$ and $D$. Then combine the results at the end.
%-------------------------------------------------------------
\subsection{Product Rule}

When the forcing term is a product of a polynomial, an
exponential, and a wave, combine their trial forms.
After expanding, give every distinct term its own
unknown coefficient.

\par\medskip
\textbf{Polynomial times exponential.}
For a forcing term such as
\[
f(t)=4te^{2t},
\]
use a full linear polynomial multiplied by the exponential:
\[
x_p(t)=(At+B)e^{2t}.
\]

\par\medskip
\textbf{Exponential times wave.}
For
\[
f(t)=e^{-3t}\sin 2t,
\]
include both sine and cosine:
\[
x_p(t)=e^{-3t}(A\cos 2t+B\sin 2t).
\]

\par\medskip
\textbf{Polynomial times wave.}
For
\[
f(t)=t\cos t,
\]
give sine and cosine their own full linear polynomials:
\[
x_p(t)=(At+B)\cos t+(Ct+D)\sin t.
\]
Do not use $(At+B)(C\cos t+D\sin t)$: that form ties
the coefficients together. The coefficients of
$t\cos t$, $\cos t$, $t\sin t$, and $\sin t$
must be independent.

%-------------------------------------------------------------

\subsection{Combining All Three Types}
For a forcing term such as $t^2e^{3t}\sin 2t$, use
\[
x_p=e^{3t}\bigl[(At^2+Bt+C)\cos 2t+(Dt^2+Et+F)\sin 2t\bigr].
\]
There are six independent coefficients. Both polynomial factors have
 degree two; the exponential rate and trigonometric frequency remain fixed.
This is a trial-selection exercise, not a request to differentiate this long expression.

% SIDE NOTE: 9:40 | Your Turn: Choose the Trial
% a) Identify the building blocks before writing coefficients.
% b) Use independent coefficients and combine repeated trial terms.
% c) Let students explain their choices; leave the long product for reading if late.
\subsection{Practice Choosing Trials}
\begin{example}{For each equation, choose a trial function. Do not determine its coefficients.
\begin{enumerate}
\item $x''-x=2t+e^{-3t}$.
\item $x''+x=e^{-3t}\sin 2t$.
\item $x''+4x=t\cos t$.
\item $x''+2x'+2x=t^2+3t+e^{2t}$.
\end{enumerate}}
\begin{enumerate}
\item $x_p=At+B+Ce^{-3t}$: add the polynomial and exponential trials.
\item $x_p=e^{-3t}(A\cos 2t+B\sin 2t)$: include both waves.
\item $x_p=(At+B)\cos t+(Ct+D)\sin t$: each wave gets its own polynomial.
\item $x_p=At^2+Bt+C+De^{2t}$: one full quadratic already covers both $t^2$ and $3t$.
Do not add another linear polynomial with redundant coefficients.
\end{enumerate}
\end{example}

\subsection{A Product Worked Out}
% OPTIONAL EXTENSION: Independent reading, or use in a longer class.
\begin{example}{Find a particular solution of $x''-x=4te^{2t}$.}
Try $x_p=(At+B)e^{2t}$. Differentiating carefully gives
\begin{align*}
x_p'&=(2At+A+2B)e^{2t},\\
x_p''&=(4At+4A+4B)e^{2t}.
\end{align*}
Substitution gives
\[
x_p''-x_p=(3At+4A+3B)e^{2t}=4te^{2t}.
\]
Cancel the nonzero factor $e^{2t}$ and match coefficients:
\[
3A=4,\qquad 4A+3B=0.
\]
Hence
\[
\boxed{x_p=\left(\frac43t-\frac{16}{9}\right)e^{2t}.}
\]
The problem asks only for a particular solution. If it asked for the general
solution, we would add $c_1e^t+c_2e^{-t}$.
\end{example}

% SIDE NOTE: 9:48 | The Main Idea
% a) Choose a family that contains the terms produced by differentiation.
% b) Substitute and match coefficients.
% c) Next class: what if no choice of coefficients works?
\subsection{Looking Ahead}
So far, substitution has allowed us to determine the coefficients in each trial.
Next time, we will encounter a problem where our usual choice does not work.
We will use the failure to understand how to repair the method.
\end{document}
